\documentclass[conference]{IEEEtran}
\usepackage[utf8]{inputenc}
\usepackage[T1]{fontenc}
\usepackage{cite}
\usepackage{amsmath,amssymb}
\usepackage{graphicx}
\usepackage{booktabs}
\usepackage{url}

\title{Verifier‑Assisted versus One‑Shot LLM Intent\ensuremath{\rightarrow}Configuration: A Paired NetOpsTraceBench Study}

\author{
\IEEEauthorblockN{Autonomous AI Researcher}
\IEEEauthorblockA{CNSM 2026 Agentic AI Researcher Track}
}

\begin{document}

\maketitle

\begin{abstract}
We preregistered and executed NetOpsTraceBench, a paired comparison of a one‑shot LLM intent\ensuremath{\rightarrow}configuration pipeline and a verifier‑assisted generate--validate--repair pipeline on 1,200 intent\ensuremath{\rightarrow}config episodes. Using gpt-5-mini, a deterministic NetOps validator, and Mininet‑style emulation, we applied a preregistered binary episode success rule (validator accepts final config AND emulator shows no availability degradation above tolerance AND no automatic rollback). Of 1,200 planned paired tasks we completed 1,199 pairs. Baseline success = 1,196/1,199 (0.9974979149); guarded success = 1,199/1,199 (1.0000000000). The paired‑success‑rate difference (guarded - baseline) = 0.0025020851 (95\% bootstrap percentile CI [0.0000, 0.0058381985], bootstrap\_seed=1729, resamples=10,000); exact McNemar two‑sided p = 0.25. The preregistered power target sought a 5 percentage‑point minimum detectable effect; given the observed near‑ceiling baseline, the study lacked power to detect much smaller absolute gains. Representative validator traces and provider audit logs show repair was invoked only three times. Under these controlled emulation and task‑corpus conditions we find no statistically significant advantage for the verifier‑assisted pipeline; we report artifacts, analytic deviation (Wilcoxon \ensuremath{\rightarrow} exact McNemar for binary outcomes), and detailed execution accounting to support reproducible interpretation.
\end{abstract}

\section{Introduction}
Generative large language models (LLMs) are increasingly proposed to translate operator intent expressed in natural language into executable device configurations for network operations (NetOps). Systems and survey literature emphasize that operational reliability depends critically on orchestration, deterministic validators, and assurance infrastructure rather than on an LLM alone [1,2]. Generate--validate--repair (verifier‑assisted) architectures are a canonical mitigation strategy: an LLM proposes candidate outputs, an automated verifier checks constraints and safety, and an automated repair loop iterates until the candidate satisfies verifier rules [3--5]. Despite conceptual support, publicly archived, preregistered, paired evidence measuring whether verifier assistance yields an operationally meaningful improvement (rollback‑aware binary success) relative to one‑shot generation is sparse.

We address this gap with NetOpsTraceBench: a preregistered, paired, emulator‑grounded study comparing (A) baseline one‑shot NL\ensuremath{\rightarrow}config generation and (B) a guarded generate--validate--repair pipeline using the same shared initial candidate. The study asked: Can a reproducible, workflow‑centered benchmark of paired intent\ensuremath{\rightarrow}config episodes produce an objective paired binary success metric that reliably distinguishes verifier‑assisted pipelines from one‑shot generation across two simulated topologies? Our contributions are (1) a preregistered experimental design and full execution artifacts (study ID rX4f7-1); (2) a paired analysis on 1,199 completed pairs using gpt-5-mini and a deterministic validator; (3) transparent reporting of an analytic deviation (final confirmatory test used an exact McNemar test appropriate for paired binary outcomes) and execution accounting that explains why the guarded pipeline rarely required repair.

\section{Related Work}
Recent surveys and system papers position LLMs within AIOps/NetOps pipelines and emphasize the importance of verifiers, provenance, and end‑to‑end workflow evaluation [1,6]. Work on verifier‑guided NL\ensuremath{\rightarrow}formal translations and intent\ensuremath{\rightarrow}topology compilation has shown the feasibility of generate--validate‑repair approaches in domains such as security topology compilation and CTL/specification generation [3,7]. Empirical and conceptual studies document LLM failure modes relevant to NetOps---hallucination, incorrect semantics/units, overconfidence, provenance loss, and prompt injection---which motivate guarded architectures and runtime auditing [8--10]. Several recent prototypes demonstrate LLM‑assisted configuration generation and misconfiguration detection but typically evaluate on constructed or small held‑out sets [4,5,11]. That literature motivates workflow‑centered evaluation metrics (trace quality, rollback measurement, canaries) rather than one‑shot QA metrics; NetOpsTraceBench operationalizes a paired, rollback‑aware binary success metric in a preregistered framework and archives the run artifacts for reproducibility.

\section{Methodology}
Preregistration and experimental scope: The study design was preregistered under ID rX4f7-1 and frozen prior to execution (preregistration manifest SHA‑256 = 77abe8b4...). The preregistration specified confirmatory hypotheses, a single primary estimand (paired\_success\_rate\_difference\_guarded\_minus\_baseline), deterministic inclusion/exclusion rules, contamination thresholds, stopping rules, random seeds, maximum model‑call budget, and a complete analysis plan. NetOpsTraceBench\_v1 targeted 1,200 paired intent\ensuremath{\rightarrow}config tasks partitioned into two simulated topology clusters (edge = 600, core = 600). Task examples were public synthetic or consented; every task carried deterministic provenance metadata. The preregistered power plan assumed a baseline success near 0.70 and set a minimum detectable effect (MDE) of 5 percentage points (absolute) at two‑sided \ensuremath{\alpha}=0.05 and power 0.8; sensitivity scenarios over a range of baseline and discordant‑pair assumptions were specified in the preregistration.

Pipelines, transformations, and instrumented components: We compared two paired conditions executed from a single shared initial candidate per pair. Baseline (one‑shot) invoked a single LLM generation call to produce an initial candidate configuration. Guarded (verifier‑assisted) used the same shared initial candidate and permitted at most one automated repair iteration when the deterministic validator rejected the candidate; preregistration limited maximum\_repair\_calls\_per\_task = 1 and initial\_generation\_calls\_per\_task = 1 to bound model‑call budget. The requested model was gpt-5-mini (provider recorded as openai\_responses) invoked through orchestration adapter hosted\_netops\_gvr\_v1. Deterministic\_netops\_task\_generator\_v5 produced canonical, vendor‑agnostic intent\ensuremath{\rightarrow}configuration tasks; deterministic\_netops\_validator\_v5 performed parsing, intent‑constraint validation, dependency checks, and emulator preflight evaluation. The execution contract capped maximum\_model\_calls = 2,400 and defined the scientific\_confirmatory execution mode.

Inclusion, contamination, and stopping rules: Inclusion required that neither episode in a pair be deterministically excluded. The preregistration implemented contamination detection by (a) a normalized substring/token‑overlap test with threshold T\_contam = 0.85 and (b) an optional embedding‑similarity secondary check (cosine threshold = 0.92) for flagged items. If either test exceeded thresholds the pair was flagged and excluded from the primary confirmatory set (but retained for exploratory reporting). Pairs were also excluded if both episodes failed to produce parsable vendor renderings after one automated re‑execution attempt. The stopping rule required running the full planned workload unless (i) cumulative model calls reached the maximum\_model\_calls cap, or (ii) automation halted due to deterministic component failure after the one allowed automated recovery attempt; all such events were logged.

Determinism, seeds, and failure handling: The preregistration froze generation seeds and analysis seeds. Deterministic task generation and validator behavior reduce sampling variance from non‑deterministic test artifacts; model nondeterminism remained possible at the provider side but was bounded by the single‑call-per‑stage policy. Missingness handling followed the preregistered plan: the primary analysis used complete\_pair\_primary (exclude flagged/excluded pairs); a sensitivity analysis treated excluded or failed episodes as failures (complete\_pair\_primary\_with\_failure\_as\_zero\_sensitivity). All exclusion reasons, failed episode counts, and provider audit traces were archived to preserve transparent accounting.

Success metric, confirmatory test, and bootstrap CI: The preregistered binary episode success rule required three conjunctive outcomes: (1) final configuration accepted by the deterministic validator (parsing and intent constraints satisfied), (2) emulator execution shows availability degradation no greater than the preregistered tolerance, and (3) no automatic rollback event occurred during simulated execution. Although the preregistration specified a paired nonparametric Wilcoxon signed‑rank test, the archived outcomes were binary for each episode. Following the preregistered emphasis on correct paired‑outcome inference, the final archived confirmatory analysis applied an exact McNemar two‑sided test to the 2\ensuremath{\times}2 paired contingency table and computed a paired bootstrap percentile 95\% confidence interval for the paired‑success‑rate difference using frozen bootstrap\_seed = 1729 and 10,000 resamples. This analytic deviation (Wilcoxon \ensuremath{\rightarrow} exact McNemar) is recorded in the analysis artifacts and in the Disclosure Statement.

Instrumentation, auditing, and artifact recording: The run captured provider call audit metadata (hosted\_model\_call\_event\_count and stage\_counts), all raw model outputs, shared initial candidate traces, repair provider traces when present, validator traces (validation\_before/validation\_after), emulator execution logs, task\_manifest entries, transformation\_manifest, and a full hash manifest. Key configured limits recorded in model\_configuration.json included max\_output\_tokens = 1500 and requested\_model = gpt-5-mini. The execution manifest records planned\_episode\_count = 2,400, completed\_episode\_count = 2,398, failed\_episode\_count = 2, and model\_calls\_used = 1,203. All seeds, analysis code (paired\_binary\_analysis\_v1), and final analysis artifacts (analysis/paired\_contingency\_table.csv, analysis/results.json, analysis/paired\_difference\_figure.svg) are archived to enable exact reproduction.

Predefined sensitivity and exploratory plans: The preregistration enumerated exploratory analyses that the archived run executed or preserved for downstream work: (a) distribution of repair iterations and their correlation with final success and time‑to‑repair; (b) automated failure‑mode categorization (semantic unit errors, missing dependencies) derived from deterministic parsers and emulator checks; (c) sensitivity of the primary binary success to alternative availability‑degradation tolerances and to treating flagged contamination pairs as failures. All exploratory outputs were labeled as such and are reported separately from the confirmatory inference to control Type I error per the preregistered multiplicity plan.

\section{Results}
Execution and accounting (archived): Planned episode\_count = 2,400 (1,200 pairs); completed\_episode\_count = 2,398; failed\_episode\_count = 2; completed\_pairs = 1,199. Provider audit: hosted\_model\_call\_event\_count = 1,203; completed\_hosted\_model\_call\_event\_count = 1,202; failed\_hosted\_model\_call\_event\_count = 1. Stage counts show shared\_initial\_generation = 1,200 and repair = 3 (i.e., the guarded pipeline triggered repair calls only three times). Artifact\_count = 8,408 and the run completed at 2026‑08‑30T03:15:12.141890+00:00; representative artifact examples and validator traces are archived (e.g., execution/scoring/task-000001-baseline.json).

Primary confirmatory outcomes (archived analysis): complete\_pair\_count = 1,199; baseline\_success\_count = 1,196 (baseline\_success\_rate = 0.9974979149291076); guarded\_success\_count = 1,199 (guarded\_success\_rate = 1.0). Paired contingency: n11 = 1,196; n10 = 3 (guarded=1 baseline=0); n01 = 0; n00 = 0 (see analysis/paired\_contingency\_table.csv). Paired‑success‑rate difference (guarded - baseline) = 0.002502085070892446. Exact McNemar two‑sided p = 0.25. Paired bootstrap percentile 95\% CI (bootstrap\_seed = 1729; 10,000 resamples) = [0.0000, 0.005838198498748957]. Missingness summary: failed\_baseline\_episodes = 1; failed\_guarded\_episodes = 1; both\_missing\_pairs = 1 (see analysis/missingness\_summary.csv).

Representative diagnostics: Validator traces in scoring artifacts show that, for typical tasks, the initial shared candidate satisfied intent constraints and validator checks (e.g., execution/scoring/task-000001-baseline.json shows validation\_after.valid = true and repair\_applied = false). The low repair invocation count (3) is corroborated by the provider audit stage\_counts. These artifact‑grounded diagnostics indicate that under the chosen validator configuration and synthetic task corpus the LLM initial candidates were accepted in the vast majority of episodes, leaving little headroom for guard‑triggered improvement.

\section{Discussion}
Statistical and operational interpretation: The exact McNemar test (p = 0.25) and the paired bootstrap CI [0.0000, 0.0058381985] both indicate the observed guarded advantage is not distinguishable from zero under our confirmatory criterion. The empirical discordant counts (n10 = 3, n01 = 0) directly determine the paired estimate: paired\_difference = 3/1,199 = 0.0025020851. Because the preregistered MDE was 0.05 (5 percentage points), the study---while sized for that target---was not powered to detect the much smaller absolute difference we observed. Concretely, observing a 5 pp absolute paired difference on the realized complete\_pair\_count (1,199) would require approximately 60 discordant pairs (0.05 \ensuremath{\times} 1,199 \ensuremath{\approx} 59.95), but we observed only three discordants; this arithmetic shows why the planned MDE and realized discordant count are inconsistent with detecting the tiny observed effect.

Artifact‑grounded diagnostics explain why repairs were rare. The provider audit and stage counts record only three repair invocations (repair = 3) and model\_calls\_used = 1,203, corroborated by the validator traces where validation\_after.valid = true and repair\_applied = false for representative tasks (see execution/scoring/task-000001-baseline.json). Two interacting design choices reduced headroom for verifier benefit: (1) the synthetic task generator produced well‑specified intents with canonical required\_sequences (deterministic\_netops\_task\_generator\_v5), and (2) the deterministic validator (deterministic\_netops\_validator\_v5) used parsing and constraint checks that accepted many reasonable initial candidates. These two factors are visible in representative task\_manifest entries (required\_sequence and difficulty fields) and in the scoring artifacts where parsed\_command\_count == required\_command\_count and violation\_count == 0 for typical episodes. Together they explain the low frequency of repair-triggered iterations and the near‑ceiling baseline success (\ensuremath{\approx}99.75\%).

Reproducibility and forensic trace evidence: All numerical claims above are reproducible from archived artifacts. The confirmatory analysis used bootstrap\_seed = 1729 with 10,000 resamples to compute the percentile CI (analysis/analysis\_log.jsonl), and the paired contingency table is published (analysis/paired\_contingency\_table.csv). Execution accounting (execution/execution\_manifest.json) reports started\_at\_utc and completed\_at\_utc, model\_calls\_used, completed pairs, and failed episodes. Artifact\_count = 8,408 and representative scoring JSON entries include validator\_trace objects showing validation\_before/after, observed\_touched\_field\_count, and repair\_applied flags. Those deterministic traces permit exact reproduction of the numeric summaries without human adjudication.

Threats to validity made concrete by artifacts: Internal validity is high because task generation, seeds, and deterministic validators were frozen in the preregistration (preregistration SHA‑256 = 77abe8b4...5712) and automation enforced exclusion rules; the missingness\_summary.csv documents the two failed episodes and the single both\_missing\_pair. Construct validity depends on how well the binary episode success rule (validator acceptance + emulator availability tolerance + no rollback) maps to operator‑valued correctness; the validator's acceptance thresholds and the emulator's availability tolerance are parameters embodied in the validator traces and transformation\_manifest and therefore materially affect measured outcomes. External validity is limited: only one model family (requested\_model = gpt-5-mini) was invoked (execution/model\_configuration.json), tasks were synthetic/consented and vendor‑agnostic, and emulation is Mininet‑style rather than vendor‑specific hardware; these constraints appear in task\_manifest entries and in the emulator preflight logs. Conclusion validity is constrained by the small number of discordant pairs: standard paired tests rely on sufficient discordance to support inference, and here the discordant mass (3/1,199) is too small to yield narrow hypothesis tests for small absolute differences.

Implications for benchmark design and operational practice: The artifacts show that when initial generation already satisfies deterministic verifier rules for well‑specified synthetic tasks, guarded architectures will seldom be exercised, and measured binary success gains will be tiny. To reveal verifier value in future studies, designers should (a) increase nominal task ambiguity or adversarial content so that initial candidates more frequently violate intent constraints, (b) vary validator strictness to control repair invocation rates, and (c) test multiple model families or fine‑tuned variants to measure model--validator interactions. The archived run provides concrete knobs: change task\_manifest difficulty fields, adjust validator violation thresholds, or expand maximum\_repair\_calls\_per\_task to study repair dynamics---each change is reproducible because seeds, manifests, and validator traces are preserved.

Operational diagnostics and sensitivity: The low repair count (3) correlated with specific task patterns (deterministic task difficulty levels labelled "medium" or "hard" in representative manifests) but was insufficient to permit reliable subgroup inference. Exploratory artifacts (analysis/condition\_summary.csv and analysis/missingness\_summary.csv) can be used to stratify by difficulty and re‑estimate discordant rates; however, such exploratory subgroup tests remain underpowered without increasing task diversity or sample size. The run's provider audit also shows a single failed hosted\_model\_call\_event\_count = 1; recovery and retry rules are logged and the stopping\_rule was enforced without manual intervention.

Concluding operational guidance: The null confirmatory result---guarded advantage not statistically significant---should be interpreted as conditional on the task corpus, validator configuration, and requested model. The run artifacts consistently point to limited headroom (near‑ceiling baseline) and low repair invocation as the proximate causes. For practitioners, the key takeaway is not that verifiers are universally unnecessary, but rather that their measured value depends strongly on task ambiguity and validator policy. We provide the full artifact set (task manifests, validator traces, provider audits, analysis logs, and seeds) to enable sensitivity experiments that vary these parameters and thereby expose conditions under which verifier scaffolding yields operationally meaningful improvements.

\section{Limitations}
\begin{itemize}
\item Single model family tested (gpt-5-mini); results do not imply generality across other LLMs or fine‑tuned variants.
\item Task corpus skew: synthetic/consented tasks yielded near‑ceiling baseline success (\ensuremath{\approx}99.75\%), limiting headroom and statistical power to detect practical improvements by guarded pipelines.
\item Emulator fidelity: Mininet‑style emulation and deterministic validators approximate production behavior but may not capture all deployment failure modes (routing convergence, vendor idiosyncrasies, transient telemetry).
\item Verifier configuration dependence: deterministic\_netops\_validator\_v5 acceptance thresholds and parsing rules materially affect outcomes; different verifier strictness could change repair invocation rates and measured gains.
\item Analytic deviation documented: the preregistration specified a Wilcoxon signed‑rank paired procedure; the archived confirmatory analysis used an exact McNemar test because outcomes were binary. This deviation is recorded in the archived analysis artifacts and in the Disclosure Statement above.
\item Operational external validity: absence of heterogeneous vendor renderers, real production templates, and live rollouts constrain direct production generalization.
\end{itemize}

\section{Conclusion}
We preregistered and executed NetOpsTraceBench, a paired, emulator‑grounded study comparing one‑shot and verifier‑assisted NL\ensuremath{\rightarrow}configuration pipelines on 1,199 completed pairs. Under our synthetic task corpus, deterministic validator settings, and gpt-5-mini as requested, baseline one‑shot generation already achieved near‑ceiling success (\ensuremath{\approx}99.75\%). The guarded pipeline increased absolute success by \ensuremath{\approx}0.25 percentage points (paired difference = 0.0025020851) with exact McNemar two‑sided p = 0.25 and a 95\% bootstrap CI that includes zero. Archived execution manifests, provider audits (model\_calls\_used = 1,203; repairs invoked = 3), representative validator traces, and analysis artifacts support the interpretation that the observed null result reflects limited headroom given task and verifier choices rather than definitive evidence that verifiers are never useful. Future work should intentionally design more challenging and adversarial task sets, vary verifier strictness, and test multiple LLM families to determine conditions under which verifier scaffolding yields operationally meaningful improvements. All run artifacts, analysis code, seeds, and logs are archived to enable exact reproduction and follow‑up sensitivity studies (see Disclosure Statement).

\section*{Disclosure Statement}
Mandatory disclosure (preregistered run rX4f7-1). LLM(s) and provider: requested model = gpt-5-mini via provider recorded as openai\_responses; model configuration archived at execution/model\_configuration.json (requested\_model=gpt-5-mini, max\_output\_tokens=1500). Orchestration/adapter: hosted\_netops\_gvr\_v1 executed the paired benchmark. Deterministic tooling: deterministic\_netops\_validator\_v5 performed canonical parsing, intent‑constraint validation, and emulator preflight checks; deterministic\_netops\_task\_generator\_v5 produced synthetic tasks. Exact initial master prompt: archived at provenance/master\_prompt.txt (immutable reference SHA‑256 = 1872df1e\allowbreak{}1805d2d9\allowbreak{}6940456c\allowbreak{}a016bd66\allowbreak{}5d1d5196\allowbreak{}add77f5a\allowbreak{}cdf1582b\allowbreak{}b39b15ba). Topic constraint and autonomy boundary: the human Research Director selected the topic family and approved the master prompt pre‑lock; after the autonomy lock the AI autonomously formulated the research question, conducted literature review, generated the preregistration, executed the experiment, performed analysis, and wrote and revised the manuscript. Preregistration: NetOpsTraceBench (ID rX4f7-1) preregistration manifest SHA‑256 = 77abe8b4\allowbreak{}93bde444\allowbreak{}84e0b3d0\allowbreak{}2cfd7f9e\allowbreak{}c1621ff7\allowbreak{}e7847889\allowbreak{}082ae177\allowbreak{}a7ed5712; preregistered design (confirmatory hypotheses, inclusion/exclusion, seeds, stopping rules, analysis plan) was frozen before execution. Execution summary: started 2026-08-30T01:21:55.415955+00:00, completed 2026-08-30T03:15:12.141890+00:00; planned episodes 2,400; completed episodes 2,398; completed pairs 1,199; model\_calls\_used = 1,203; repair-stage calls observed = 3. Analysis artifacts and deviation record: the preregistration specified a Wilcoxon signed‑rank paired procedure; the archived executed confirmatory analysis used an exact McNemar two‑sided test on paired binary outcomes plus a paired bootstrap percentile CI. This post‑preregistration analytic choice is documented in the archived analysis artifacts (analysis/paired\_contingency\_table.csv, analysis/results.json, analysis/analysis\_log.jsonl, analysis/paired\_difference\_figure.svg) produced as part of the run that completed at 2026-08-30T03:15:12.141890+00:00. Artifacts and reproducibility: raw model outputs, provider traces, validator traces, task\_manifest, transformation\_manifest, execution logs, analysis code, and all seeds are archived; artifact\_count = 8,408 and full hash manifest path = execution/execution\_manifest.json. Human involvement: pre‑lock collaboration and infrastructure development occurred; no human manual editing, selective revision, or ad hoc text insertion occurred on the manuscript content after the autonomy lock---the AI performed internal autonomous revision cycles only. Technical restarts and deviations: execution manifest records two failed episodes and the analytic deviation described above; all deviations are logged in execution/execution\_log.jsonl. The disclosure above is complete relative to the archived artifacts supplied with this run.

\begin{thebibliography}{99}
\bibitem{ref1} http://hdl.handle.net/20.500.12708/229494
\bibitem{ref2} https://arxiv.org/abs/2608.13389
\bibitem{ref3} http://arxiv.org/abs/2403.09369
\bibitem{ref4} https://doi.org/10.2139/ssrn.6947162
\bibitem{ref5} https://doi.org/10.2139/ssrn.7222278
\bibitem{ref6} https://doi.org/10.20935/acadai8260
\bibitem{ref7} https://doi.org/10.3390/info17030297
\bibitem{ref8} https://doi.org/10.3390/app16010085
\end{thebibliography}

\end{document}
